\subsection{Definition und Gradvergleich}
\label{subsec:real_functions_rational_fundamentals}
\par\noindent\begin{minipage}{\linewidth}
\begin{DefinitionBox}{Gebrochen-rationale Funktion}{real_functions_rational}
  Eine Funktion $f$ heißt \emph{gebrochen-rationale Funktion}
  \EnglishTerm{rational function}, wenn sie sich als Quotient zweier Polynome
  schreiben lässt:
  \[
    f(x)=\frac{Z(x)}{N(x)}
      =\frac{a_mx^m+a_{m-1}x^{m-1}+\dots+a_0}
            {b_nx^n+b_{n-1}x^{n-1}+\dots+b_0},
    \qquad
    \begin{aligned}
      \deg(Z)&=m,\\
      \deg(N)&=n.
    \end{aligned}
  \]

  \medskip
  Die Nullstellen des Nenners müssen aus dem Definitionsbereich ausgeschlossen werden:
  \[
    D_f=\{x\in\R\mid N(x)\ne0\}.
  \]
\end{DefinitionBox}

\begin{DefinitionBox}{Echt und unecht gebrochene Funktionen}{real_functions_proper_improper_rational}
  Eine gebrochen-rationale Funktion $f(x)=Z(x)/N(x)$ mit $Z\not\equiv0$ heißt
  \begin{itemize}
    \item \emph{echt gebrochen} \EnglishTerm{proper rational function}, wenn $\deg(Z)<\deg(N)$;
    \item \emph{unecht gebrochen} \EnglishTerm{improper rational function}, wenn $\deg(Z)\ge\deg(N)$.
  \end{itemize}
\end{DefinitionBox}

% Handschrift: PDF-Seite 23 (alle drei Gradvergleichsbeispiele).
\begin{ExampleBox}{Zählergrad und Nennergrad vergleichen}{real_functions_rational_classification}
  Ordnen Sie die Funktionen anhand der höchsten vorkommenden Potenzen ein.
  Die Definitionsmenge wird weiterhin aus dem ursprünglichen Nenner bestimmt.
  \begin{enumerate}[label=\alph*),itemsep=0.8\baselineskip]
    \item \(\displaystyle
        \mathmakebox[3.2cm][r]{f(x)=\frac{x^2+x+1}{x^4-1}}\quad\Longrightarrow\quad
        \begin{aligned}
          \deg(Z)&=\Cloze{0.8cm}{2}\\
          \deg(N)&=\Cloze{0.8cm}{4}
        \end{aligned}\quad\Longrightarrow\quad
        \text{\Cloze{4cm}{echt gebrochen}}
      \)
    \item \(\displaystyle
        \mathmakebox[3.2cm][r]{g(x)=\frac{x^4-1}{x^3+1}}\quad\Longrightarrow\quad
        \begin{aligned}
          \deg(Z)&=\Cloze{0.8cm}{4}\\
          \deg(N)&=\Cloze{0.8cm}{3}
        \end{aligned}\quad\Longrightarrow\quad
        \text{\Cloze{4cm}{unecht gebrochen}}
      \)
    \item \(\displaystyle
        \mathmakebox[3.2cm][r]{h(x)=\frac{x^4+1}{x^4-1}}\quad\Longrightarrow\quad
        \begin{aligned}
          \deg(Z)&=\Cloze{0.8cm}{4}\\
          \deg(N)&=\Cloze{0.8cm}{4}
        \end{aligned}\quad\Longrightarrow\quad
        \text{\Cloze{4cm}{unecht gebrochen}}
      \)
  \end{enumerate}
\end{ExampleBox}
\end{minipage}
